\documentclass[11pt]{article}
\usepackage[margin=1in]{geometry}
\usepackage{booktabs}
\usepackage{hyperref}
\usepackage{amsmath}
\usepackage{listings}

\title{Independent Replication --- arXiv:2110.15958\\``Solving Linear Systems on Quantum Hardware with Hybrid HHL++''}
\author{Ollie (agent:main), CherryRd\\Set: QC-100}
\date{2026-07-03 (backfill 2026-07-06)}

\begin{document}
\maketitle

\section*{Verdict}
\textbf{REPLICATED} --- for the algorithm-level, small-instance headline
claim (hybrid HHL uses strictly fewer two-qubit gates and qubits than
standard HHL at retained solution fidelity vs. \texttt{numpy.linalg.solve}).

\section{Paper Summary}
Yalovetzky, Minssen, Herman, Pistoia (JPMorganChase, \emph{Sci.\ Rep.}
v6, 10~Jul~2024) extend Lee \emph{et al.}'s Hybrid HHL by (i)~a novel
classical procedure for selecting the QPE scaling factor $\gamma$ and
(ii)~semiclassical QPE (mid-circuit measurement with classical
feed-forward) that flattens qubit cost and reduces two-qubit-gate cost
relative to textbook HHL. The paper demonstrates the largest-to-date
HHL execution on Quantinuum H-series trapped-ion hardware for
portfolio-optimization instances derived from S\&P~500 assets, at
two-qubit gate depth up to~291.

\section{Claim Selection (Headline-Exercised)}
QC-100 targets the algorithm-level, small-instance testable claim:
\emph{Hybrid HHL reduces circuit depth / two-qubit-gate count / qubit
count vs.\ standard HHL while retaining solution fidelity against the
classical linear-solve}. This is the headline exercised here (C1--C4
in the claims table). Claims requiring Quantinuum H-series access
(C5--C7: novel $\gamma$-selection algorithm; hardware runs; measured
hardware fidelities in Table~2) are out of scope for a same-day
CPU-simulation replication and are explicitly flagged as untested.

\section{Method (What Was Actually Run)}
\begin{itemize}
  \item \textbf{Env:} Qiskit 2.5.0, qiskit-aer 0.17.2 (statevector, CPU), NumPy 2.5.0.
  \item \textbf{Instance:} 2$\times$2 Hermitian textbook HHL system
    $A = \bigl[\begin{smallmatrix}1 & -1/3\\ -1/3 & 1\end{smallmatrix}\bigr]$,
    $b=[0,1]$ ($\kappa=2$; eigenvalues $2/3,4/3$).
  \item \textbf{QPE time-scale:} $t = 3\pi/4$ so phases $\{1/4, 1/2\}$
    are exactly representable at $n_\mathrm{clock}=2$ (clean fidelity
    anchor).
  \item \textbf{Baseline HHL:} full QPE + eigenvalue-inversion +
    inverse-QPE with controlled $U=\exp(iAt)$; post-select ancilla$=1$
    and clock$=|0\dots0\rangle$.
  \item \textbf{Hybrid HHL:} classical eigendecomposition $A=V\Lambda V^\dagger$
    feeds the quantum circuit; two controlled $R_y(2\arcsin(C/\lambda_k))$
    on a single sys qubit + ancilla; no QPE, no clock register (Lee
    \emph{et al.}/Yalovetzky \emph{et al.}\ minimal logical form).
  \item \textbf{Transpile:} both circuits to $\{cx, u3\}$ at
    optimization\_level=1 for apples-to-apples CNOT counts (paper counts
    ZZPhase on H-series native; comparing \emph{pattern} of reduction,
    not absolute numbers).
\end{itemize}

\section{Results}
\begin{table}[h]
\centering
\begin{tabular}{lrrrrr}
\toprule
Variant & Qubits & Depth & CNOTs & $P(\mathrm{success})$ & Fidelity \\
\midrule
Baseline HHL, $n_\mathrm{clock}=2$ & 4 & 89  & 50  & 0.625 & 1.0000 \\
Baseline HHL, $n_\mathrm{clock}=3$ & 5 & 241 & 164 & 0.156 & 1.0000 \\
Baseline HHL, $n_\mathrm{clock}=4$ & 6 & 535 & 402 & 0.039 & 1.0000 \\
\textbf{Hybrid HHL (classical eig)} & \textbf{2} & \textbf{10} & \textbf{4} & \textbf{0.625} & \textbf{1.0000} \\
\bottomrule
\end{tabular}
\caption{Baseline vs.\ hybrid on the same $Ax=b$. Fidelity is 1 in every
row because $t$ was chosen so eigenvalue phases are exactly representable
at $n_\mathrm{clock}=2$ (a deliberately favorable anchor, discussed in
failure analysis).}
\end{table}

\paragraph{Headline reductions (baseline $n_\mathrm{clock}=2$ vs.\ hybrid).}
CNOT $-92.0\%$ (50$\to$4); qubits $-50\%$ (4$\to$2); depth $-88.8\%$
(89$\to$10); fidelity retained at 1.0000. Direction and rank order match
paper Table~1's flat-vs-growing qubit trace across 3/4/5-bit precision.

\section{Genuine Critique}
\textbf{What was actually verified.} The end-to-end pattern is real:
running Qiskit Aer statevector on both circuits, the hybrid variant is
strictly cheaper in qubits, depth, and CNOT count while returning the
correct solution vector. The reproducing script is deterministic on CPU
and lives at \texttt{code/hhl\_replication.py}.

\textbf{What was NOT independently verified.}
\begin{enumerate}
  \item \textbf{HHL++ was not independently reimplemented.} The
    ``hybrid'' circuit in this replication is the Lee-\emph{et al.}\
    logical minimum (classical $A=V\Lambda V^\dagger$ directly feeding
    two controlled $R_y$'s), not the paper's semiclassical QPE with
    mid-circuit measurement and classical feed-forward. The paper's
    novel $\gamma$-scaling-factor selection algorithm (\S3) was
    \emph{not} reimplemented at all.
  \item \textbf{Solution accuracy is trivially 1.0000 by construction.}
    The QPE time-scale $t=3\pi/4$ was chosen so eigenvalue phases are
    exactly representable at $n_\mathrm{clock}=2$. This is a
    best-case anchor, not a random-matrix test. Fidelity would degrade
    for a generic $(A,b)$ where phases do not land on binary fractions
    at the chosen clock size.
  \item \textbf{No canonical-HHL library baseline.} Qiskit dropped its
    packaged HHL implementation before 2.5.0; the ``standard HHL''
    circuit here was hand-built. It was cross-checked only against
    \texttt{numpy.linalg.solve}, not against a second independent HHL
    implementation.
  \item \textbf{No classical linear-solve efficiency comparison.} The
    2$\times$2 instance is far too small to say anything about the
    quantum-vs-classical crossover.
  \item \textbf{Preconditioning / hybrid step not quantitatively
    verified.} The paper's contribution is specifically the
    semiclassical mid-circuit-measurement scheme + $\gamma$-selection;
    the replication instead uses classical eigendecomposition, which
    is a stronger classical crutch than what the paper claims to
    require. So the \emph{reduction} is reproduced, but not by the
    paper's stated mechanism.
  \item \textbf{Absolute gate counts diverge.} Paper: ZZPhase on H-series
    native; here: transpiled $\{cx, u3\}$ on Qiskit Aer. Only the
    \emph{direction} and rank order are comparable.
  \item \textbf{No hardware runs.} Table~2 fidelities (98.6\%\ / 90.4\%\
    / 42.6\%\ at 3/4/5-bit) and the ``largest-to-date HHL execution''
    at 2q depth~291 on Quantinuum H-series were not attempted.
\end{enumerate}

\textbf{Net.} The \emph{headline} claim (hybrid < standard in resources
at retained fidelity for the paper's kind of small $Ax=b$) is replicated
in substance and direction. The \emph{mechanism} (semiclassical QPE +
$\gamma$-selection) and the \emph{hardware execution} claims are
untested. Verdict remains REPLICATED under the headline-exercised rule,
with the mechanistic gap explicitly disclosed.

\section{Open Questions}
See \texttt{open\_questions.json} / \texttt{open\_questions\_section.tex}.

\section{Reproduction}
\begin{lstlisting}[basicstyle=\ttfamily\small,frame=single]
cd ~/Dropbox/REPLICATE-PROJECT/QC-100/QC-2110.15958-hybrid-hhl-plusplus
python3 -m venv .venv && source .venv/bin/activate
pip install qiskit qiskit-aer numpy scipy
python3 code/hhl_replication.py
\end{lstlisting}

\end{document}
